\caption{\rev{\textbf{Synergistic causality and its hierarchical organization in controlled dynamical systems.}
\textbf{a}, Factorized uniform interventions remove source dependence while keeping the transition mechanism fixed; under the adopted PID axioms, the redundancy term vanishes.
\textbf{b}, Schematic of a nonlinear system with a fixed $\sin(xy)$ interaction and a common driver $w$, whose influence is controlled by $\beta$.
\textbf{c}, Response of MLP+PEID (teal solid line), WMS, MMI-PID, SURD and MLP+SHAP to increasing common-driver strength $\beta$.
\textbf{d}, Generating topologies, SPTs and hierarchical clustering outputs for three six-node benchmarks. I, mixed pairwise/triadic Kuramoto; II, Glauber dynamics on the same interaction support; III, Kuramoto with $\{1,2,3\}\subset\{1,2,3,4\}$ and module $56$, connected by two weak pairwise bridges. Green, red and grey borders denote planted modules, unsupported merges and unscored refinements or roots, respectively. SPT numbers and teal fills show local Syn as a percentage of the tree's total, using one common scale.
\textbf{e}, Module $F_1$ for the displayed seed-0 outputs. Singleton leaves, the common root and refinements inside primitive modules are excluded from scoring. MIC (data) uses samples directly. III reports each baseline's better of two predefined data views.}}
\label{fig:synergy_overview}
